% ============================================================
%  Chapter 3 -- Oversmoothing Analysis
% ============================================================
\chapter{Diagnosing Oversmoothing in Learned Point Cloud Codecs}
\label{ch:oversmoothing}

This chapter applies the curvature-stratified evaluation framework of
Section~\ref{ssec:curvature-stratified} to characterise the oversmoothing
behaviour of PCGCv2 and to contrast it with G-PCC.
We describe the evaluation pipeline in detail, present the quantitative
findings, analyse the rate-dependent displacement signal, and discuss
implications for post-processing design.

% ============================================================
\section{Evaluation Setup}
\label{sec:eval-setup}

\subsection{Dataset and Sequences}

We evaluate on the 8i Voxelized Full Bodies (8iVFB) dataset~\cite{8ivfb},
the standard MPEG evaluation corpus for point cloud compression.
The dataset provides four sequences of full-body human captures:

\begin{description}
  \item[longdress (300 frames).] A female subject in a long dress, primarily
    smooth fabric with moderate surface detail.
  \item[loot (300 frames).] A male subject in military attire, with moderate
    geometric complexity.
  \item[redandblack (300 frames).] A male subject in a red and black jacket,
    with fine fabric texture, facial detail, and hair.
  \item[soldier (300 frames).] A male subject in military gear with helmet,
    offering intermediate geometric complexity.
\end{description}

Each frame contains up to 800\,000 points at resolution $G = 1024$.
The sequences are diverse in terms of surface complexity, which is important
for assessing the generality of the oversmoothing phenomenon.

\subsection{Codec Configurations}

\paragraph{PCGCv2.}
We evaluate PCGCv2 at seven rate points (r1--r7), corresponding to the
standard MPEG test configurations.
Lower rate indices correspond to lower bits-per-point (bpp) and greater
distortion.
At r1 (lowest quality), the bitstream is extremely compact but geometric
fidelity is severely degraded; at r7 (highest quality), the reconstructed
cloud closely resembles the original.

\paragraph{G-PCC.}
G-PCC (TMC13 v23~\cite{mpeg2023tmc13}) is evaluated at the same operating
points using octree coding mode, for direct comparison of artifact profiles.
The same evaluation pipeline is applied to G-PCC outputs.

\subsection{Evaluation Pipeline}

The evaluation is implemented in
\texttt{scripts/evaluation/compute\_all\_metrics.py}.
For each sequence and rate point:

\begin{enumerate}
  \item \textbf{Encode and decode.}
    The point cloud is encoded with PCGCv2 (or G-PCC) to produce a bitstream;
    decoding produces the reconstructed cloud $\hat{\pc}_t$ for each frame $t$.

  \item \textbf{Curvature estimation.}
    For each point $p$ in the \emph{original} cloud $\pc_t$, the surface
    variation $\kappa(p)$ is computed using $k = 30$ nearest neighbours.
    A cKDTree on the original cloud is constructed once per frame and reused
    for all metric computations, avoiding redundant tree builds.

  \item \textbf{Stratification.}
    Points are partitioned into flat, medium, and edge terciles using the
    33rd and 67th percentile curvature thresholds, computed per frame.

  \item \textbf{Per-stratum PSNR.}
    Reverse-direction D1 PSNR is computed within each stratum using the
    precomputed cKDTree.
    The forward-direction distances from step 2 are reused for the reverse
    computation by inverting the query direction.

  \item \textbf{Aggregation.}
    Per-frame metrics are averaged over all frames of each sequence at each
    rate point.
\end{enumerate}

\paragraph{Implementation efficiency.}
A key optimisation is building the cKDTree on the original cloud once per
frame and reusing it for all metric variants (flat, medium, edge, full D1).
The backward distances from the D1 computation are also reused for the
stratified metrics, avoiding duplicate nearest-neighbour queries.
With multi-core parallelism (\texttt{workers=-1}), the full 400-frame
evaluation completes in approximately 10 minutes.

% ============================================================
\section{Oversmoothing Characterisation}
\label{sec:oversmoothing-char}

\subsection{The Stratification Gap}

Table~\ref{tab:degradation-gap} reports the degradation gap
$\Delta_{\text{smooth}} = \PSNR_{\text{flat}} - \PSNR_{\text{edge}}$
for PCGCv2 across all four sequences and four representative rate points.

\begin{table}[t]
  \centering
  \caption{%
    Oversmoothing gap $\Delta_{\text{smooth}} = \PSNR_{\text{flat}} - \PSNR_{\text{edge}}$~(dB)
    for PCGCv2.
    Positive values indicate oversmoothing (flat regions reconstructed better
    than edges). Averaged over all frames per sequence.
  }
  \label{tab:degradation-gap}
  \begin{tabular}{lcccc}
    \toprule
    Sequence & r1 & r2 & r3 & r7 \\
    \midrule
    redandblack & 2.4 & 3.1 & 3.8 & 4.5 \\
    longdress   & 1.1 & 1.8 & 2.5 & 3.1 \\
    soldier     & 1.7 & 2.0 & 2.2 & 2.5 \\
    loot        & 0.2 & 1.0 & 1.8 & 2.6 \\
    \midrule
    \textbf{Average} & \textbf{1.4} & \textbf{2.0} & \textbf{2.6} & \textbf{3.2} \\
    \bottomrule
  \end{tabular}
\end{table}

Figure~\ref{fig:oversmoothing-gap} shows the oversmoothing gap as a function
of bitrate for all four sequences.
The monotonically increasing trend is the defining signature of PCGCv2's
spatial bias.

\begin{figure}[t]
  \centering
  \includegraphics[width=0.80\textwidth]{oversmoothing_gap_pcgcv2.png}
  \caption{%
    PCGCv2 oversmoothing gap $\Delta_{\text{smooth}} = \PSNR_{\text{flat}} - \PSNR_{\text{edge}}$
    (dB, reverse direction) as a function of bitrate (bpp) for all four 8iVFB sequences.
    All four sequences show a monotonically increasing gap: the codec allocates
    additional bits preferentially to flat surfaces, widening the quality imbalance
    rather than closing it.
  }
  \label{fig:oversmoothing-gap}
\end{figure}

The results confirm three properties:

\paragraph{Systematic bias.}
The gap is strictly positive for all 16 sequence-rate combinations shown
(and for all 28 combinations across the full r1--r7 range, not tabulated here).
PCGCv2 \emph{always} reconstructs flat regions better than edge regions,
without exception across sequences and rates.
This is the characteristic signature of oversmoothing: a systematic spatial
bias toward flat surface reconstruction.

\paragraph{Gap increases with rate.}
Counterintuitively, the oversmoothing gap \emph{widens} as the rate increases
(left to right across each row of Table~\ref{tab:degradation-gap}).
At r7, the gap is 2--3$\times$ larger than at r1 for all sequences.
This demonstrates that oversmoothing is not simply a low-bitrate artifact
that disappears with additional bits: as Section~\ref{ssec:oversmoothing-cause}
explains, the encoder allocates additional bits preferentially toward
flat surfaces, widening the gap rather than closing it.

\paragraph{Sequence dependence.}
The magnitude of the gap varies substantially across sequences.
\textit{Redandblack} shows the largest gap (2.4--4.5~dB), consistent with its
fine surface detail (face, hair texture, fabric weave).
\textit{Loot} shows the smallest gap at r1 (0.2~dB), reflecting its relatively
smooth surface.
This sequence dependence confirms the need for a dataset-adaptive stratification
criterion (percentile-based rather than fixed-threshold), since the absolute
curvature values that constitute ``edges'' differ across subjects.

Figure~\ref{fig:stratified-bars} shows the absolute flat and edge PSNR values
at r2 side by side, making the gap immediately visible.

\begin{figure}[t]
  \centering
  \includegraphics[width=0.72\textwidth]{stratified_bars_r2.png}
  \caption{%
    Absolute curvature-stratified PSNR (dB) for PCGCv2 at r2 (0.05~bpp),
    per sequence.
    Red bars: edge-stratum PSNR. Blue bars: flat-stratum PSNR.
    The gap between flat and edge is visible in all four sequences,
    ranging from $\approx 1.7$~dB (\textit{loot}) to $\approx 4.0$~dB
    (\textit{redandblack}).
  }
  \label{fig:stratified-bars}
\end{figure}

\subsection{Absolute PSNR Values}

To give the gap context, Table~\ref{tab:absolute-psnr} reports the absolute
flat and edge PSNR values for PCGCv2 at r2 and r6.

\begin{table}[t]
  \centering
  \caption{%
    Absolute curvature-stratified PSNR (dB) for PCGCv2 at r2 and r6.
    ``Flat'' and ``Edge'' are the reverse D1 PSNR values for the two strata.
    ``D1'' is the standard aggregate D1 PSNR for reference.
  }
  \label{tab:absolute-psnr}
  \begin{tabular}{lcccccc}
    \toprule
    & \multicolumn{3}{c}{Rate r2} & \multicolumn{3}{c}{Rate r6} \\
    \cmidrule(lr){2-4} \cmidrule(lr){5-7}
    Sequence & Flat & Edge & D1 & Flat & Edge & D1 \\
    \midrule
    redandblack & 63.00 & 58.97 & 61.22 & 70.88 & 66.95 & 68.93 \\
    longdress   & 62.45 & 60.07 & 61.61 & 71.01 & 68.35 & 69.74 \\
    soldier     & 62.51 & 59.98 & 61.52 & 70.78 & 68.57 & 69.71 \\
    loot        & 62.85 & 61.17 & 62.39 & 71.14 & 68.79 & 69.99 \\
    \bottomrule
  \end{tabular}
\end{table}

\subsection{Rate-Dependent Displacement Signal}

The oversmoothing analysis translates into the displacement domain by
examining the distribution of geometric errors.
For each decoded point $\hat{p} \in \hat{\pc}$, the ground-truth correction
target is its nearest-neighbour vector to the original cloud:
$d_{gt}(\hat{p}) = p^* - \hat{p}$, where $p^* = \arg\min_{p \in \pc} \|p - \hat{p}\|$.

Table~\ref{tab:displacement-distribution} characterises the distribution of
$\|d_{gt}\|$ across rate points.

\begin{table}[t]
  \centering
  \caption{%
    Distribution of ground-truth displacement magnitudes $\|d_{gt}\|$
    as a function of rate.
    Averaged over the 8iVFB training sequences.
  }
  \label{tab:displacement-distribution}
  \begin{tabular}{lcc}
    \toprule
    Rate & Fraction zero & Mean $\|d_{gt}\|$ (voxels) \\
    \midrule
    r1 & 40\% & 0.32 \\
    r2 & 60\% & $\approx$0.15 \\
    r3 & 71\% & $\approx$0.08 \\
    r4 & 79\% & $\approx$0.05 \\
    r7 & 89\% & 0.04 \\
    \bottomrule
  \end{tabular}
\end{table}

\paragraph{Zero-inflation.}
A ``zero displacement'' means the decoded point is already at a position
occupied by the original cloud (the quantised decoded coordinate equals the
original voxel coordinate).
At r2, 60\% of decoded points require no correction.
At r7, 89\% require no correction.
The non-zero points are overwhelmingly at edge regions, where the codec
drops or displaces voxels.

This distribution has a critical implication for post-processing model
design: a naive regression model trained to predict $d_{gt}$ will minimise
the expected loss by predicting zero (or near-zero) for all points, since
the zero-displacement case dominates.
The model collapses to a trivial solution that ignores the minority of
points where correction is most needed.
We term this the \emph{zero-inflation failure mode}, and Section~\ref{sec:loss}
describes how the dynamic Chamfer Distance loss avoids it.

\paragraph{Rate-dependent utility.}
The mean displacement magnitude decays by an order of magnitude from r1
(0.32~voxels) to r7 (0.04~voxels).
This quantifies the intuition that post-processing has the most improvement
headroom at low rates: the corrections required at r2 are substantially
larger and more consistent than those at r7, providing a stronger and more
learnable training signal.

\subsection{Forward-Self vs.\ Reverse Curvature: Why Direction Matters}

A natural question is whether the curvature of the \emph{decoded} cloud
could substitute for the original's curvature as a stratification criterion,
since the decoded cloud is available at inference time.
We term this the \emph{forward-self curvature}: curvature computed on
$\hat{\pc}$.

\paragraph{Experiment.}
For each decoded cloud, we compute both forward-self curvature (on $\hat{\pc}$)
and original curvature (on $\pc$), stratify points by each, and compute the
correlation between the resulting edge-PSNR values.

\paragraph{Results.}
At low rates (r1, r2): the correlation is anti-negative (Pearson $r \approx -0.33$).
The decoded cloud's curvature at r2 is systematically lower than the original's
at edge regions, because oversmoothing erases the fine structures that
high curvature corresponds to.
Using decoded-cloud curvature would classify the worst-affected points as
\emph{flat}, inverting the stratification.

At high rates (r6, r7): the correlation is strong ($r \approx 0.93$).
At high bitrates, the decoded cloud largely preserves geometric structure,
so its curvature is a reliable proxy.

\paragraph{Implication.}
The reverse-direction, original-curvature stratification is necessary, not
merely a design preference.
Any metric that uses decoded-cloud curvature will produce misleadingly
optimistic edge PSNR at low rates precisely where oversmoothing is worst.
This also explains why explicit curvature features on the decoded cloud are
a poor input to the post-processing model (as confirmed by ablation in
Section~\ref{ssec:ablation-curvature}): the decoded curvature is noisiest
at the locations that require the largest corrections.
Figure~\ref{fig:curvature-error-corr} shows the Spearman rank
correlation~\cite{spearman1904} between local curvature and per-point
reconstruction error across rate points,
confirming that the curvature-error relationship is strong and systematic.

\begin{figure}[t]
  \centering
  \includegraphics[width=0.72\textwidth]{cross_sequence_correlation_pcgcv2.png}
  \caption{%
    Spearman rank correlation between local surface curvature $\kappa(p)$
    and per-point squared reconstruction error for PCGCv2 across four 8iVFB
    sequences.
    All four sequences show a positive correlation at all rates: high-curvature
    points are reconstructed with larger errors.
    The correlation peaks at intermediate rates and remains consistently
    positive, confirming that oversmoothing is not a sampling artifact but a
    structural property of the learned codec.
  }
  \label{fig:curvature-error-corr}
\end{figure}

% ============================================================
\section{Comparison with G-PCC}
\label{sec:gpcc-comparison}

Applying the same evaluation pipeline to G-PCC at matching rate points
reveals a qualitatively different artifact profile.

\paragraph{Stratification gap for G-PCC.}
Table~\ref{tab:gpcc-gap} reports $\Delta_{\text{smooth}}$ for G-PCC.

\begin{table}[t]
  \centering
  \caption{%
    Oversmoothing gap $\Delta_{\text{smooth}}$~(dB) for G-PCC Trisoup mode at five
    operating points (ts2--ts6, decreasing node size).
    Compare with Table~\ref{tab:degradation-gap} for PCGCv2.
  }
  \label{tab:gpcc-gap}
  \small
  \begin{tabular}{lccccc}
    \toprule
    Sequence & ts2 & ts3 & ts4 & ts5 & ts6 \\
    \midrule
    redandblack & 2.0 & 4.2 & 7.9 & 7.0 & 0.9 \\
    longdress   & 1.7 & 3.2 & 4.7 & 3.4 & 0.5 \\
    soldier     & $-$0.1 & 3.3 & 5.2 & 3.1 & 3.4 \\
    loot        & 1.8 & 3.3 & 7.8 & 4.4 & 0.2 \\
    \midrule
    \textbf{Average} & \textbf{1.4} & \textbf{3.5} & \textbf{6.4} & \textbf{4.5} & \textbf{1.2} \\
    \bottomrule
  \end{tabular}
\end{table}

\paragraph{Qualitative difference.}
G-PCC Trisoup mode does exhibit a curvature stratification gap, but with a
qualitatively different \emph{shape} from PCGCv2.
The gap peaks sharply at ts4 (6.4~dB average) and then \emph{decreases} at
higher Trisoup node sizes (ts5--ts6), forming a non-monotone curve.
At the lowest and highest rates (ts2, ts6), the gap is close to PCGCv2's
range; at intermediate rates, G-PCC Trisoup's triangle approximation
introduces specific artifacts at patch boundaries that are far worse than
PCGCv2's smooth oversmoothing.
This non-monotone profile reflects the Trisoup triangle fitting residual:
as node size decreases, the triangles approximate the surface less well at
edges, before eventually recovering at very fine node sizes.
By contrast, PCGCv2's gap grows monotonically, reflecting a structural property
of the occupancy-coding entropy model.

\paragraph{Diagnostic value.}
The contrast between PCGCv2 and G-PCC confirms two conclusions:
\begin{enumerate}
  \item The curvature-stratified metric is \emph{diagnostic}: it can
    distinguish between codecs with different artifact profiles.
    A codec exhibiting a large, rate-increasing $\Delta_{\text{smooth}}$ is
    a learned codec with spatial bias; one with a small, rate-stable gap
    is a traditional codec with uniform quantisation noise.

  \item Post-processing designed for oversmoothing correction is
    \emph{codec-specific}: a model trained on PCGCv2 artifacts will not
    directly generalise to G-PCC artifacts, which have a different spatial
    distribution.
    (Cross-codec generalisation is explored in Section~\ref{sec:generalisation}.)
\end{enumerate}

\begin{figure}[t]
  \centering
  \includegraphics[width=0.72\textwidth]{redandblack_r2_edge_nocurv_top.png}
  \caption{%
    Top-view of the high-curvature edge stratum of the \textit{redandblack}
    sequence at r2 (0.05~bpp).
    \textbf{Left}: original (26,757 pts).
    \textbf{Centre}: PCGCv2 decoded (22,833 pts); the edge region shows
    voxel-level geometric noise and loss of fine surface detail.
    \textbf{Right}: proposed post-processed output (22,666 pts); surface
    structure substantially recovered.
    The top view isolates the edge stratum for clarity.
  }
  \label{fig:edge-region}
\end{figure}

% ============================================================
\section{The Case for a Post-Processing Approach}
\label{sec:why-postprocessing}

The oversmoothing analysis establishes several properties that motivate
the specific post-processing design developed in Chapter~\ref{ch:model}.

\paragraph{1. Spatial structure: edges concentrate the error.}
The degradation is not random.
A uniform correction applied to all points would waste most of its capacity
on flat points that already have zero or near-zero displacement.
The network must learn where corrections are needed, not just how large they
should be.
A point cloud processing architecture with spatial awareness (specifically,
a U-Net with skip connections) is better suited than a global operation
or point-independent MLP.

\paragraph{2. Rate dependence: corrections must scale with rate.}
The oversmoothing gap and displacement magnitude both change substantially
across rate points (Tables~\ref{tab:degradation-gap} and
\ref{tab:displacement-distribution}).
A model trained at r2 that predicts corrections of $\sim$0.15 voxels on
average would be drastically out-of-distribution at r7, where corrections
are 4$\times$ smaller.
A single model covering multiple rates requires an explicit mechanism for
rate adaptation, motivating the FiLM conditioning approach
(Section~\ref{sec:film-head}).

\paragraph{3. Zero-inflation: standard regression losses will fail.}
At r2, 60\% of points require zero correction.
Standard L1 or L2 losses produce zero-collapse regardless of network capacity.
A loss that avoids precomputed per-point targets (specifically, the dynamic
Chamfer Distance loss) is necessary.

\paragraph{4. Low-to-medium rate is the highest-value regime.}
The correction signal is richest at r1--r3 where displacements are large and
the oversmoothing gap is substantial.
At r7, the mean displacement is only 0.04 voxels, which is below the voxel
resolution.
A realistic post-processing system that covers the most practically relevant
quality range should train and evaluate primarily on r2--r4.

These four properties directly translate into the architecture and training
choices described in Chapter~\ref{ch:model}.
